\section{Results and Analysis}
\label{sec:results}

\subsection{Main Results}

Table~\ref{tab:scorecard} gives the final state of all six runs. The central empirical result is that PPO and GRPO respond to variance filtering in opposite directions.

\begin{table*}[t]
\centering
\small
\begin{tabular}{lllllll}
\toprule
\# & Run & Algorithm & Filter & Reward & Success & Format \\
\midrule
1 & PPO vanilla & PPO & 1.0 & \(-0.654\) & 0.000 & 0.000 \\
2 & PPO StarPO-S moderate & PPO & 0.5 & \(-0.185\) & 0.000 & 0.000 \\
3 & PPO StarPO-S aggressive & PPO & 0.25 & \(-0.060\) & 0.005 & 0.050 \\
4 & GRPO vanilla & GRPO & 1.0 & \(\mathbf{+0.225}\) & \(\mathbf{0.225}\) & \(\mathbf{0.808}\) \\
5 & GRPO filtered moderate & GRPO & 0.5 & \(-0.079\) & 0.005 & 0.369 \\
6 & GRPO filtered aggressive & GRPO & 0.25 & \(-0.072\) & 0.020 & 0.363 \\
\bottomrule
\end{tabular}
\caption{Final evaluation metrics at step 200. GRPO vanilla is the only run with sustained positive reward, non-trivial success, and high format compliance.}
\label{tab:scorecard}
\end{table*}

\begin{figure*}[t]
\centering
\includegraphics[width=\figw]{figures/fig01_eval_reward_6groups.png}
\caption{Evaluation reward over all six runs. PPO follows the expected filtering trade-off direction, while GRPO performs best without variance filtering.}
\label{fig:reward}
\end{figure*}

\paragraph{PPO directionally reproduces RAGEN.}
PPO with filter ratio 1.0 deteriorates from \(-0.124\) at step 20 to \(-0.654\) at step 200, reproducing the vanilla instability direction in the RAGEN small-model sparse-reward setting. Moderate filtering delays and reduces the collapse: at step 120, PPO with ratio 0.5 reaches \(-0.054\), while vanilla PPO is already at \(-0.516\). At the final step, the moderate filter remains 0.469 reward points better than vanilla. Aggressive filtering is stable but nearly stagnant, so its final score should not be read as a clean algorithmic win.

\paragraph{GRPO reverses the StarPO-S pattern.}
GRPO vanilla is the strongest configuration in the entire matrix, reaching reward \(+0.225\), success rate 0.225, and format compliance 0.808. Filtering does not repair GRPO; both filtered GRPO variants stay near the negative-reward plateau. This is the main extension beyond reproduction: the variance filter that helps PPO does not transfer directly to actor-only GRPO.

Supplementary success, format, and multi-metric curves are moved to Appendix~\ref{sec:app_figs}. They support the same conclusion: GRPO vanilla is strongest or tied strongest across the main evaluation dimensions, while filtered GRPO is consistently weaker.

\subsection{Why Filtering Helps PPO but Hurts GRPO}

To explain why the same variance filter helps PPO but hurts GRPO, I separate the mechanism into three parts: whether high-variance groups improve critic learning, how costly sample removal is for each algorithm, and when aggressive filtering degenerates into a single-mini-batch update.

\paragraph{Critic-data selection exists only in PPO.}
For PPO, high-variance prompts provide useful value-function targets. The same prompt can produce trajectories with different returns, so the critic receives informative state-return contrasts. Better critic targets improve advantage estimates and can stabilize the actor update. Variance filtering therefore acts partly as data selection for the critic.

GRPO has no critic. Its advantage is a within-group z-score. Once a group has nonzero reward variation, scaling the absolute reward spread does not proportionally scale the learning signal. Filtering high-variance groups therefore does not give GRPO the same signal amplification; it mostly removes samples.

\paragraph{Sample removal is more costly in GRPO.}
For PPO, cleaner critic data can partially offset the cost of fewer trajectories, producing a practical sweet spot. For GRPO, this benefit is much weaker. Fewer samples increase gradient variance and reduce the number of update mini-batches without a compensating critic-training gain.

\paragraph{The 0.25 boundary is shared.}
With \(P=8\), \(K=16\), and effective mini-batch size 32, the number of optimizer steps per training iteration is
\begin{equation}
n_{\mathrm{grad}} =
\left\lceil \frac{r \cdot P \cdot K}{32} \right\rceil.
\end{equation}
Thus ratios 1.0, 0.5, and 0.25 produce 4, 2, and 1 optimizer steps. At ratio 0.25, the filtered buffer contains exactly one mini-batch. With one epoch, the first and only update has current log probabilities equal to old log probabilities, so the PPO-style clipping mechanism never activates. This degeneration affects both PPO and GRPO, explaining why the aggressive-filter runs look stable but stagnant.

\subsection{Failure Mode Diagnostics}

The vanilla PPO failure is not the same as RAGEN's classic echo trap. RAGEN echo trap is usually associated with repeated invalid modes and decreasing entropy. Here, PPO vanilla shows falling reward, zero format compliance, a sharp KL spike, and increasing entropy. This is better described as format collapse: the policy moves into malformed high-entropy output rather than a deterministic repeated phrase.

GRPO vanilla has the opposite shape. Its entropy decreases while reward, success, and format compliance increase. This is healthy mode convergence, not echo trap. The diagnostic lesson is that entropy alone is ambiguous; reward and format compliance must be interpreted jointly. Appendix~\ref{sec:app_figs} includes the supporting entropy, KL/gradient, head-to-head, and echo-proxy figures.

\subsection{Quantization and Algorithm Scale}

The PPO runs have pre-clip gradient norms in the hundreds or thousands, while GRPO remains around \(0.05\)--\(0.16\). This matches the different advantage scales: PPO's bi-level GAE can produce large token-level weights, while GRPO z-scores are naturally bounded. Since AdamW8bit introduces block-wise optimizer-state quantization noise~\citep{dettmers2022optimizers}, PPO's larger update scale and critic coupling may be more sensitive to numerical drift. GRPO provides a useful contrast: it uses the same hardware concessions but avoids the critic and uses smaller advantages, and it is the only configuration that learns robustly in this setting.
